%% kapa: exam=1 round=29 year=2018 period=1 form=A subject=civil_law subject_ko=민법 seq=36 number=36
%% source: Q-Net 감정평가사 기출문제 (cst00309) artlSeq=5202808 file=extracted/1008/5202808/2018년 제29회 감정평가사 1차시험 기출문제/2018년 제29회 감정평가사 1차시험 기출문제/1교시 A형.pdf
%% conversion: kordoc 4.24.51 -> tools/gen1.py (auto); review=proofread-sonnet; flags=none
\begin{question}{36}
공동소유에 관한 설명으로 옳지 않은 것은? (다툼이 있으면 판례에 따름)
\begin{choices}
  \item 비법인사단이 타인 간의 금전채무를 보증하는 행위는 총유물의 관리ㆍ처분행위이므로 사원총회의 결의를 요한다.
  \item 토지 공유자의 공유지분 포기에 따른 등기는 해당 지분에 관하여 다른 공유자 앞으로 소유권이전등기를 하는 형태가 되어야 한다.
  \item 합유물에 관하여 경료된 무효의 소유권이전등기 말소청구는 특별한 사정이 없는 한, 합유자 각자가 할 수 있다.
  \item 공유물분할의 소는 공유자 전원이 소송당사자로 참여해야 하므로, 공동소송인 중 1인에 소송요건의 흠결이 있는 경우 전 소송이 부적법하게 된다.
  \item 과반수지분권자로부터 공유토지의 특정부분의 점유를 허락받은 제3자는 소수지분권자에 대해서 그 점유로 인한 이득을 부당이득으로 반환할 필요가 없다.
\end{choices}
\end{question}
\answer{1}
